\section{Discussion}
% Budget: discussion, limitations, and statements share 0.9 page.
% Order: answer each measurement question from supported results; explain correction
% failure or success; explain chain-only synthetic consequences; state limitations;
% close with the evidence needed before any use with people.
% Define lack of fit and conditional simulation uncertainty in reader-facing terms.
% Leave out project history, health associations, deployment advice, and human claims.
\pending{Interpret each feature decision and correction result within the tested range. Explain synthetic consequences under their stated assumptions.}

\subsection{Limitations}
% Write one limitations paragraph. Cover one tested unit of each device; software,
% surfaces, mountings, angle, azimuth, load, active region, sizes, and speeds.
% Explain that a robot cannot measure changes in human movement caused by surfaces.
% Cite R6 for the surface effect. The correction addresses only measurement-chain error.
% Band power concerns designed line oscillations only. Jerk, pressure, tilt, and hover
% are descriptive. The camera is not a band-power or instantaneous-timing reference.
% Explain conservative feature decisions and limited power without treating failure
% as proof of a difference. Explain that omitted-mounting checks assess lack of fit.
% Synthetic Gaussian feature assumptions, local-shift transfer, population scaling,
% independent joint features, and fixed bench estimates constrain interpretation.
% Simulation intervals exclude bench-estimation uncertainty. Behavioral and chain
% effects can oppose each other. Do not state a numerical lower bound for people.
\pending{Write the limitations paragraph. Explain the surface effect \cite{R6} and each bound on inference.}
